\begin{center}\begin{tikzpicture}[font=\small,>=stealth,line join=round]
\draw[gray!40,rounded corners=3pt] (0,0) rectangle (14.2,8.3);\node at (7.1,7.700000000000001) {图像证明：相对于中间水平线，两次上下换边};
\begin{scope}[shift={(1.7000000000000002,2.139361702127659)},x=1.2cm,y=0.978723404255319cm]\draw[->,gray] (-0.5,0)--(4.5,0) node[right]{$x$};\draw[->,gray] (0,-0.5)--(0,4.2) node[above]{$y$};\draw[blue,thick,solid,domain=-0.2:4.2,samples=180] plot(\x,{2.5+cos(deg(pi*\x))});\fill[orange!85!black] (0,3.5) circle(2.2pt);\fill[orange!85!black] (1,1.5) circle(2.2pt);\fill[orange!85!black] (2,3.5) circle(2.2pt);\draw[violet,dashed] (-.4,2.5)--(4.4,2.5);\end{scope}\node[anchor=west,align=left,text width=5.0cm] at (8.4,4.0) {示例 $f=\frac52+\cos(\pi x)$\\[9pt]中间高度 $\frac52$\\[9pt]$f(x+1)=5-f(x)$\\[9pt]两次换边，回到原高度\\[9pt]$f(x+2)=f(x)$};
\node[text width=12.6cm,align=center] at (7.1,.65) {虚线是两次高度的平均位置；换边与横移同步，不能把它误作周期1。};
\end{tikzpicture}\end{center}
